\paragraph{Notation.}
For a positive integer $r$, write $[r]:=\{1,\ldots,r\}$.  For a
prime $q$, let $\F_q$ denote the finite field with $q$ elements,
and let $\F_q^\times:=\F_q\setminus\{0\}$ denote its multiplicative group.
Vectors and matrices are represented in bold, e.g., $\mathbf v$ for a vector and $\mathbf A$ for a matrix, while their scalar coordinates are not. We use $\|\cdot\|_2$ for the $\ell_2$-norm and $\mathbf1_r$ for the vector of ones in $\R^r$.

\paragraph{Hasse derivatives.}
We will use Hasse derivatives in \cref{sec:descent} to recognize when a
polynomial has a zero of prescribed order at $X=1$.  Let $F$ be a field, let
$f(X)=\sum_{i=0}^d c_iX^i\in F[X]$, and let $j\ge0$ be an integer.  The $j$th
\emph{Hasse derivative} of $f$ is
\[
 D^{[j]}f(X):=\sum_{i=j}^d\binom ij c_iX^{i-j},
\]
where each integer binomial coefficient is read as an element of $F$.
We also use the binomial polynomials $\binom{X}{0}:=1$ and
$\binom{X}{j}:=X(X-1)\cdots(X-j+1)/j!\in F[X]$ for $j\ge1$ whenever
$j!\ne0$ in $F$.
For $b\in F$ and an integer $k\ge1$, we say that $b$ is a zero of order at
least $k$ when $(X-b)^k$ divides $f(X)$.
The following standard characterization will be
used~\cite{LidlNiederreiter97}:
\begin{equation}
\label{eq:hasse-multiplicity}
 (X-b)^k\mid f(X)
 \quad\Longleftrightarrow\quad
 D^{[j]}f(b)=0\quad(0\le j<k).
\end{equation}

\subsection{Euclidean lattices}
A Euclidean lattice of rank $n$ is a discrete additive subgroup
$\mathcal L\subseteq\R^m$ of the form
\[
 \mathcal L=\mathcal L(\mathbf B):=\mathbf B\Z^n,
\]
where $\mathbf B\in\R^{m\times n}$ has linearly independent columns.  These
columns form a basis.
The first minimum of $\mathcal L$ is the length of a shortest
nonzero lattice vector:
\[
 \lambda_1(\mathcal L)
 =\min_{\mathbf0\ne\mathbf v\in\mathcal L}\|\mathbf v\|_2.
\]

For computational problems, the basis
$\mathbf B\in\Q^{m\times n}$ is given explicitly.  Exact decision-$\SVP$
also receives a nonnegative rational squared threshold $\Delta$ and asks
whether
\[
 \lambda_1(\mathcal L(\mathbf B))^2\le\Delta.
\]
The problem belongs to $\NP$: a YES certificate is a nonzero integer
coefficient vector for a short lattice point.  Standard cofactor bounds give
a certificate of polynomial bit length, and exact rational arithmetic checks
both membership and the squared norm. See, e.g.,~\cite{MicciancioGoldwasser02}.
The corresponding search-$\SVP$ problem takes a full-column-rank
basis $\mathbf B$ and outputs a shortest nonzero vector of
$\mathcal L(\mathbf B)$.

\subsection{Prime cyclotomic fields and module lattices}

\paragraph{Cyclotomic field.}
Fix an odd prime $q$ for the remainder of the construction, and let
$\zeta:=\zeta_q\in\C$ be a primitive $q$th root of unity.  We suppress the
subscript $q$ after fixing the prime and put
\[
 K:=\Q(\zeta),
 \qquad
 \mathcal{O}_K:=\Z[\zeta],
 \qquad
 \pi:=1-\zeta.
\]
Here $\mathcal{O}_K$ is the ring of integers of $K$.  The elements
$1,\zeta,\ldots,\zeta^{q-2}$ form an integral basis of $K$, so
$[K:\Q]=q-1$.  The prime $q$ is totally ramified in $\mathcal{O}_K$, and hence
$(q)=(\pi)^{q-1}$~\cite{Washington97}.  Whenever $a\in\F_q$ occurs as an
exponent of $\zeta$, we identify it with its unique integer representative in
$\{0,\ldots,q-1\}$.  A representation
$v=\sum_{a\in\F_q}c_a\zeta^a$ with $c_a\in\Q$ is called a
\emph{$q$-coordinate expression}
for $v$.  Its $q$ coefficients are indexed by $\F_q$ and need not be unique,
because $\sum_{a\in\F_q}\zeta^a=0$.  By contrast, the expression of $v$ in
the power basis $1,\zeta,\ldots,\zeta^{q-2}$ is unique.

\paragraph{Canonical embedding.}
For $v\in K$, write $\overline v$ for its complex conjugate.  In particular,
$\overline\zeta=\zeta^{-1}$.  For $a\in\F_q^\times$, let
$\sigma_a:K\hookrightarrow\C$ be the embedding defined by
$\sigma_a(\zeta)=\zeta^a$.  These are all the embeddings of $K$.  The field
trace is
\[
 \Tr_{K/\Q}(v):=\sum_{a\in\F_q^\times}\sigma_a(v).
\]
 The canonical embedding sends $v$ to
 $(\sigma_a(v))_{a\in\F_q^\times}$.  For each integer $s\ge1$, we
 equip $K^s$ with the real inner product
\[
 \langle\mathbf x,\mathbf y\rangle
 :=\operatorname{Re}\sum_{i=1}^s\sum_{a\in\F_q^\times}
   \sigma_a(x_i)\overline{\sigma_a(y_i)}
\]
and its induced product norm
\[
 \|\mathbf x\|^2:=\langle\mathbf x,\mathbf x\rangle
 =\sum_{i=1}^s\Tr_{K/\Q}(x_i\overline{x_i}).
\]
For $s=1$, this is the canonical norm on $K$:
\[
 \|v\|^2=\sum_{a\in\F_q^\times}|\sigma_a(v)|^2
 =\Tr_{K/\Q}(v\overline v).
\]
The roots of unity in $K$ are exactly $\pm\zeta^j$ for $0\le j<q$
~\cite{Washington97}.  If $\varepsilon$ is a root of unity, then
$\|\varepsilon\mathbf x-\varepsilon\mathbf y\|=\|\mathbf x-\mathbf y\|$
for all $\mathbf x,\mathbf y\in K^s$.  Thus multiplication by
$\varepsilon$ is an isometry.  In particular, it preserves distances.

\paragraph{Module lattices.}
For an integer $s\ge1$, an \emph{$\mathcal{O}_K$-module lattice} in
$K^s$ is a finitely generated $\mathcal{O}_K$-submodule
$\mathcal M\subseteq K^s$, equipped with the componentwise canonical
embedding and the product norm.  It has \emph{full rank} if its $K$-span is
$K^s$.  When $s=1$, a nonzero ideal
$\mathfrak a\subseteq\mathcal{O}_K$ gives a module lattice in $K$ called an
\emph{ideal lattice}.  Both module and ideal lattices are discrete additive
subgroups of their real spans.

In this paper, we only use free module lattices of full rank.  If
$\mathbf m_1,\ldots,\mathbf m_s\in K^s$ are linearly independent over $K$, then
\[
 \mathcal M=\mathcal{O}_K\mathbf m_1+\cdots+\mathcal{O}_K\mathbf m_s
\]
is free of module rank $s$, has full rank, and has $\Z$-rank
$s(q-1)$.  The square matrix with
columns $\mathbf m_1,\ldots,\mathbf m_s$ has nonzero determinant precisely
when these vectors are linearly independent over $K$.  If
$\mathcal M\subseteq\mathcal{O}_K^s$, then full rank is equivalent
to finite index in $\mathcal{O}_K^s$.

The first minimum of an ideal lattice or module lattice is the length of its
shortest nonzero vector in the canonical or product norm.  For $v\in K$ and
an ideal lattice $\mathfrak a$, define
\[
 \dist(v,\mathfrak a):=\min_{y\in\mathfrak a}\|v-y\|.
\]

\paragraph{Coefficient representatives.}
For $\mathbf z=(z_a)_{a\in\F_q}\in\Q^q$, write
\[
 [\mathbf z]:=\sum_{a\in\F_q}z_a\zeta^a\in K.
\]
If
$\mathbf z\in\Z^q$, then $[\mathbf z]\in\mathcal{O}_K$, and we call
$\mathbf z$ an \emph{integral coefficient representative} of
$[\mathbf z]$.  Since $\Phi_q(X)=1+X+\cdots+X^{q-1}$ is the minimal
polynomial of $\zeta$, two integral coefficient vectors represent the same
element precisely when they differ by an integer multiple of $\mathbf1_q$.

The following identity converts the canonical inner product and norm into
calculations with these $q$ coefficients.
\begin{lemma}
\label{lem:cyclotomic-form}
For all rational coefficient vectors $\mathbf x,\mathbf y\in\Q^q$,
\[
 \langle[\mathbf x],[\mathbf y]\rangle
 =q\sum_a x_ay_a-
 \left(\sum_a x_a\right)\left(\sum_a y_a\right).
\]
In particular, for every $\mathbf z\in\Z^q$,
\[
 \|[\mathbf z]\|^2
 =q\sum_a z_a^2-\left(\sum_a z_a\right)^2
 =\sum_{a<b}(z_a-z_b)^2.
\]
\end{lemma}

\begin{proof}
For $a,b\in\F_q$, additive-character orthogonality gives
\[
 \sum_{t\in\F_q^\times}\zeta^{t(a-b)}
 =\begin{cases}
   q-1,&a=b,\\
   -1,&a\ne b.
  \end{cases}
\]
Hence, for rational coefficient vectors $\mathbf x$ and $\mathbf y$,
\[
 \langle[\mathbf x],[\mathbf y]\rangle
 =q\sum_a x_ay_a-\sum_{a,b}x_ay_b
 =q\sum_a x_ay_a-\left(\sum_a x_a\right)\left(\sum_b y_b\right).
\]
Taking $\mathbf x=\mathbf y=\mathbf z$ gives the norm identity.  Expanding
$\sum_{a<b}(z_a-z_b)^2$ gives its final form.
\end{proof}

\subsection{X3C and Sidon sequences}

An Exact Cover by 3-Sets (X3C) instance consists of a universe
$[M]$ and three-element sets $C_1,\ldots,C_n\subseteq[M]$.  Its incidence
matrix $\mathbf A\in\{0,1\}^{M\times n}$ is defined by $A_{ij}=1$ when
$i\in C_j$ and $A_{ij}=0$ otherwise.  A \emph{selection vector}
$\boldsymbol\xi\in\{0,1\}^n$ specifies a choice of sets, with $\xi_j=1$
exactly when $C_j$ is selected.  We call $\boldsymbol\xi$ an
\emph{X3C witness} if
\[
 \mathbf A\boldsymbol\xi=\mathbf1_M.
\]
This equality says that every universe element belongs to exactly one
selected set.  The instance is positive exactly when it has an X3C witness.
X3C is $\NP$-complete~\cite{GareyJohnson79}.

A universe element is \emph{isolated} if it belongs to none of the sets $C_j$.  An instance with an isolated element is negative.  After handling
this case separately, every row of $\mathbf A$ is nonzero and $M\le3n$.
Adjoining a disjoint three-element block together with its unique set
preserves the answer, so we may also assume $n\ge1$.

To index the $n$ selection bits, we choose distinct nonnegative
integers whose nonzero ordered differences are all distinct.  Such a sequence
is called a \emph{Sidon sequence}.  The following lemma gives a
deterministic construction with largest entry $O(n^2)$.
\begin{lemma}
\label{lem:sidon}
For every $n\ge1$, one can deterministically compute integers
\[
 0\le\alpha_1<\cdots<\alpha_n=: \alpha_\star,
 \qquad \alpha_\star=O(n^2),
\]
such that all nonzero ordered differences
$\alpha_\ell-\alpha_j$, $j\ne\ell$, are distinct.
\end{lemma}

\begin{proof}
Use the deterministic quadratic-residue Sidon construction of
Erd\H{o}s and Tur\'an~\cite{ErdosTuran41}.  For an explicit construction and
the verification of the directed-difference property, see also
Liu--Feng--Pan~\cite{LiuFengPan26}.
\end{proof}

\subsection{Power-sum estimates}

Let $q$ be prime, let $2\le k<m$, and let
$\mathbf b=(b_1,\ldots,b_{k-1})\in\F_q^{k-1}$.  Define the
\emph{power-sum solution set}
\[
 \mathcal X_{k,m}(\mathbf b)
 :=\left\{\mathbf z=(z_1,\ldots,z_m)\in\F_q^m:
   \sum_{r=1}^m z_r^j=b_j\quad(1\le j<k)\right\}.
\]
The following proposition records the uniform point-count
bounds used below.

\begin{proposition}[{\cite{wan2026np}}]
\label{prop:wan-power-sum-estimates}
Let $q$ be prime and suppose that
$2\le k\le m-2<q-2$.  For every
$\mathbf b\in\F_q^{k-1}$, the following estimates hold.
\begin{enumerate}[label=(\roman*)]
 \item
 \[
  \left||\mathcal X_{k,m}(\mathbf b)|-q^{m-k+1}\right|
  \le \frac12(2k)^m q^{(m-k+2)/2}.
 \]
 \item For distinct $i,\ell\in[m]$,
 \[
  \left|
   \bigl|\{\mathbf z\in\mathcal X_{k,m}(\mathbf b):z_i=z_\ell\}\bigr|
   -q^{m-k}
  \right|
  \le \frac12(2k)^{m-1}q^{(m-k+2)/2}.
 \]
 \item For $i\in[m]$ and $a\in\F_q$,
 \[
  \left|
   \bigl|\{\mathbf z\in\mathcal X_{k,m}(\mathbf b):z_i=a\}\bigr|
   -q^{m-k}
  \right|
  \le \frac12(2k)^{m-1}q^{(m-k+1)/2}.
 \]
\end{enumerate}
All three estimates are uniform in $\mathbf b$, and the last two are also
uniform in the displayed choices of indices and field elements.
\end{proposition}

\begin{proof}
Wan~\cite[Propositions~5.2 and~6.2]{wan2026np} proves
part~(i) and proves part~(ii) for $z_1=z_2$, respectively.  The statements in
that paper use one particular vector of prescribed power sums $\mathbf b$.
The displayed estimates and their proofs depend only on $q$, $k$, and $m$,
so they hold uniformly for every $\mathbf b$.

For distinct $i,\ell\in[m]$, a permutation of the $m$ coordinates sends the
condition $z_i=z_\ell$ to $z_1=z_2$ and leaves each equation
$\sum_{r=1}^m z_r^j=b_j$ unchanged.  Therefore Wan's estimate for
$z_1=z_2$ proves part~(ii) for every pair $i\ne\ell$.

For part~(iii), fix $z_i=a$ and delete that coordinate.  The remaining
$m-1$ variables satisfy $\sum_{r\ne i}z_r^j=b_j-a^j$ for $1\le j<k$.
Applying Wan's first point-count estimate to these $m-1$ variables proves
part~(iii).
\end{proof}
